\section{Discussion}
\label{sec:discussion}

\subsection{What the method changes}

The central effect of \AFD\ is not the addition of one more file detector. It
changes the object on which delivery authority is based. A conventional
pipeline often scans the source object and, if no blocking rule fires, releases
that same object. The security decision is then coupled to the coverage and
interpretation of the detectors. In the method developed here, exactly three
source-evidence classes select a constrained reconstruction path. They are the
normalized portable name, caller-supplied media type, and bounded byte evidence.
Delivery authority applies to the newly produced object only after an
independent parser has established its complete output profile. This change is
modest at the level of application flow, yet
substantial at the level of the claim. The accepted object is no longer merely
an object for which no detector raised an alarm. It is an object for which a
specified conjunction of positive structural conditions holds.

Positive conditions do not make detection obsolete. Signature rules remain
useful for excluding known indicators before and after reconstruction.
Unsupported active classes remain useful policy categories. Operating-system
malware protection and application isolation remain necessary. The composition
changes the role of these controls. A signature result is one rejection layer
rather than the definition of format validity. A MIME recognizer is one source
of routing evidence rather than an authority. A codec success is evidence that
a representation could be decoded rather than evidence that the complete
output belongs to the intended canonical grammar.

This separation also improves the interpretation of negative results. When a
sample is blocked because its final suffix conflicts with its complete byte
profile, the evidence supports a classifier-disagreement claim. When a sample
is blocked because bytes follow the accepted terminator, the evidence supports
a complete-consumption claim. When a source metadata element is absent from a
reconstructed and reparsed output, the evidence supports a named metadata
exclusion. None of these outcomes needs the broader assertion that a malware
family was detected.

\subsection{Canonicality as a policy-relative property}

Canonicality is sometimes used to imply a unique encoding for an abstract
object. That interpretation would be too strong for several media families in
this study. Lossy encoders may produce different byte strings for perceptually
equivalent media. EBML allows more than one size-width representation in some
positions. Container brands and codec configurations can admit alternatives
that are valid but not interchangeable under one deployment policy. The method
therefore defines canonicality relative to a concrete format profile and a
versioned policy.

This policy-relative definition retains the properties needed at the delivery
boundary. The output parser consumes the complete object. Components appear
only at admitted locations. Required components have bounded multiplicities and
consistent cross-references. Forbidden metadata is absent. Type evidence
converges. Resource constraints hold. A different policy may define a different
accepted language, but that language is explicit and pinned to the result.

The definition also avoids a hidden dependence on encoder behavior. A handler
cannot claim canonical output merely because it invoked a named encoder. The
encoder output is reparsed under the same profile that governs delivery. This
independent gate detects trailing data, unexpected application records,
container features introduced by a tool version, and handler mistakes that a
successful semantic decode would not reveal. The independence is architectural,
not absolute. A handler and validator can share a specification error.
Differential validation with separately implemented tools is therefore retained
as empirical evidence rather than treated as a proof.

\subsection{Structural and semantic reconstruction}

The distinction between structural and semantic reconstruction is the most
important limit on the disarm claim. Structural reconstruction can provide
strong container-level results. A parser can reject invalid lengths and
offsets, remove unadmitted boxes, replace comment packets, reconstruct checksums,
normalize metadata, and establish that no bytes remain beyond the declared object.
These properties directly address many prefix, suffix, nested-object, and
privacy-metadata scenarios.

Structural reconstruction can nevertheless preserve attacker-controlled codec
payload. An accepted APNG or animated-WebP frame payload, FLAC frame, MP3
frame, Ogg media packet, ISO BMFF sample, or WebM block remains an input to a codec. Deep structural validation narrows
the syntax and can verify checksums and cross-field consistency. It does not
prove that an external decoder contains no vulnerability triggered by the
retained payload. Any statement about neutralizing content inside that payload
therefore requires semantic decode-and-reencode or a decision to block.

Semantic reconstruction provides a stronger information-flow break between
source container bytes and output container bytes. Static raster images and GIF
frames are reconstructed from decoded pixel state. A high-risk audio or video
profile can require semantic transcoding. This stronger path introduces its own
trusted computing base. The decoder consumes hostile input before the new
representation exists. Process isolation, resource limits, protocol
restrictions, cancellation, and output validation reduce exposure, but do not
turn an unverified codec into verified code. The correct interpretation is
therefore conditional: given successful bounded semantic processing within the
stated isolation assumptions, source container and compressed payload bytes are
not copied into the accepted output.

The semantic-or-block policy makes the distinction operational. If a required
semantic path is absent or fails, the system does not reinterpret structural
success as an acceptable substitute. This behavior may increase false blocking
and latency. It is still preferable to silently weakening the property that
motivated the policy.

\subsection{Closest work and residual novelty}

The closest work narrows the novelty claim materially. Open Image CDR already
provides semantic JPEG reconstruction with pixel manipulation and usability
measurement~\cite{belkind2023}. ImSan already rebuilds images to remove
metadata and extraneous content in image-based polyglot cases~\cite{koch2025}.
PdfCDR already defines and evaluates rule-informed reconstruction for
PDF~\cite{dubin2023}. SafeDocs already pursues verified parsers and
high-assurance translation into safe format subsets~\cite{darpaSafeDocs}. The
present work therefore does not claim invention of CDR, decode-and-reencode,
metadata stripping, image-polyglot sanitization, safe subsetting, or verified
parser research.

The residual novelty is the following composition. Reconstruction is
endpoint-local and user-benefiting within an end-to-end encrypted workflow,
while the relay remains blind and receives no content-derived report. Exactly
three routing sources select the candidate profile, followed by an independent
complete output parser that participates in delivery authorization. The
required reconstruction strength is indexed by policy and cannot be silently
weakened. Publication is clean-only and independently validated. A
pre-specified evidence contract fixes claims, strata, denominators, and release
gates before empirical values are admitted. This is a system-boundary and
evidence contribution, not a priority claim over the underlying reconstruction
or parser techniques.

\subsection{Implications for malware-carrier threats}

The taxonomy in \cref{sec:malware-taxonomy} separates five dimensions that are
often collapsed in security reporting: replication mechanism, initial-access
mechanism, operational role, payload objective, and attachment-carrier
property. This separation explains both the relevance and the limit of \AFD\
for virus, worm, Trojan, downloader, dropper, stager, ransomware, and spyware
scenarios.

A file-infector virus may create a carrier with an executable host grammar, an
appended body, or altered structural boundaries. Code labelled as a worm must
self-propagate after execution, although its first carrier may be a script,
executable, shortcut, archive, or active document. A downloader obtains a
follow-on payload from another source. A dropper places an embedded payload. A
stager prepares or transfers control to a later stage. A Trojan may exploit a
misleading name or declared media type. Ransomware and spyware may use any of
these delivery routes before displaying their characteristic behavior. The
evaluated media-only policy blocks unsupported active classes and tests
measurable classifier, boundary, nested-object, and true-polyglot properties.
It does not execute the payload and infer replication, operational role, or
objective.

This approach may appear conservative because it refuses an aggregate
\emph{malware blocked} rate. The alternative would be scientifically weaker.
Such a rate would combine unsupported executables, malformed images, removed
privacy metadata, decoder regressions, and semantic transcodes under one
denominator.
The mechanisms and failure conditions would be hidden. Reporting by carrier
property permits reviewers to determine which threats were actually exercised
and which output condition was established.

There are also threats for which reconstruction is the wrong method. A
phishing instruction rendered as ordinary pixels remains ordinary image
content. A steganographic payload can remain in reconstructed pixels or
samples. A visually benign QR code can direct a user to a hostile location.
An audio waveform can carry commands for another recognizer. Addressing these
cases requires perceptual, semantic, provenance, or behavioral analysis with a
different error model. Incorporating such analysis into the same product may
be useful, but its conclusions must remain separate from canonical
reconstruction.

\subsection{Privacy and deployment}

Endpoint placement follows from the confidentiality model rather than from an
assumption that endpoints are always safe. The sender necessarily has the
source plaintext before encryption. The receiver necessarily obtains plaintext
after decryption if the attachment is to be used. These are the points at which
content-derived reconstruction can occur without granting the relay a new
capability. The method consequently preserves the relay's blindness while
reducing the set of representations exposed beyond the local attachment
boundary.

This endpoint placement is client-side defense for the participant's benefit,
not centrally targeted and reporting-oriented client-side scanning. The latter
searches cleartext for designated targets and may disclose a match or its
source to an outside party, which creates the distinct risks analyzed by
Abelson et al.~\cite{abelson2024}. \AFD{} neither receives a central target list
for content matching nor sends a content-derived report to the relay or an
authority. The shared use of endpoint plaintext does not erase this difference
in purpose, information flow, and control.

The same local placement creates deployment obligations. Policy updates,
format profiles, decoder versions, and signature rules become endpoint software
state. Different endpoint versions can reach different decisions. The system
must therefore record the policy and implementation version with each decision
and define compatibility behavior explicitly. A conversation participant
should not infer that another participant applied the same local profile unless
the messaging protocol authenticates such an assertion. This manuscript does
not add that assertion to the transport protocol.

Privacy leakage is reduced for metadata fields that are structurally excluded
from the published object. EXIF location, application names, comments, and
selected container records can be removed under the corresponding profile.
Passive presence of those fields is not spyware. Their absence is a field-level
statement, not complete de-identification. Identifying information may be
visible in pixels, audible in samples, inferable from content, or encoded in a
retained stream. The evidence tables therefore report each excluded metadata
class rather than a general privacy score.

\subsection{Policy selection and operational cost}

One universal profile is unlikely to serve every deployment. A media-only
profile with semantic reconstruction for all admitted formats offers the
strongest source-to-output separation but can reject formats for which no
pinned semantic path exists. A structural profile can preserve more media and
reduce transcoding cost while supporting only container and metadata claims. A
low-latency profile may use tighter size and duration limits. A child-safety or
high-assurance deployment may reject animation, video, or all formats whose
semantic path relies on a large external decoder.

The policy algebra provides a disciplined way to compare such profiles, but
not every pair is ordered. One profile may accept PNG and reject video, while
another accepts structurally reconstructed video but imposes stricter image
dimensions. Neither is globally stronger without a declared comparison domain.
Monotonicity tests are therefore applied only to pairs for which an order has
been stated. This prevents the word \emph{strict} from substituting for a formal
relation.

Performance must be interpreted in the same profile-specific manner.
Reconstructing a small static raster in process has a different cost structure
from starting an external video transcode. Rejecting a malformed prefix may be
much cheaper than parsing a valid long-duration stream. Median latency alone
can conceal these distinctions. The evaluation stratifies by format,
reconstruction level, size, decoded size, and outcome, and retains upper
quantiles, memory, expansion, and sequential service rate. Usability is
measured independently through decoding completion and fidelity. A fast result
that blocks most benign inputs is not a successful system result.

\subsection{Evidence hierarchy}

The manuscript distinguishes analytical argument, implementation inspection,
automated property evidence, fuzzing evidence, corpus evidence, and performance
measurement. These categories answer different questions. The acceptance
predicate supports an analytical argument that clean-only publication follows
if all conjuncts and interface invariants are enforced. Source inspection
connects those conjuncts to implementation sites. Property tests seek
counterexamples over generated inputs. Fuzzing searches parser state spaces
under a recorded budget. Corpus evaluation measures behavior over declared
strata. Performance experiments estimate costs in one pinned environment.

No category silently promotes itself into another. Fuzzing without a crash does
not prove parser safety. A bounded model does not verify an external codec.
Agreement between two parsers does not prove the specification unambiguous. A
proof argument about verdict composition does not establish that every handler
correctly implements a media standard. A clean Rust build does not prove
memory safety across foreign libraries and child processes. This hierarchy is
not rhetorical caution. It defines what evidence a future reviewer would need
to challenge each claim.

At the present drafting stage, the analytical construction and implementation
traceability can be inspected. The pre-specified evidence contract fixes the
empirical protocol and its result macros remain pending. Effectiveness,
false-blocking, fidelity, and performance conclusions are not available until
the controlled server run satisfies the release gates. The absence of
provisional numbers is intentional.

\subsection{Limitations and threats to validity}

The first construct-validity threat is the mapping from campaign terminology to
carrier properties. Public labels can be incomplete or inconsistent, and one
malware family can use many carriers. The study mitigates this threat by
evaluating the property directly and retaining the behavioral label only as
scenario context. It does not infer family coverage from a synthetic carrier.

The second threat is corpus representativeness. A finite corpus cannot span all
legal encodings, malformed structures, true polyglots, trailing-object
constructions, nested objects, or decoder vulnerabilities. Boundary generators,
mutations, regression cases,
differential parsers, and coverage-guided fuzzing broaden the search. They do
not establish population-level absence of defects. Results are reported by
stratum and profile so that corpus composition remains visible.

The third threat is ground truth. Extensions and upload labels are inadequate.
The corpus contract therefore requires reproducible generation, agreement of
complete reference parsers, or documented expert adjudication. Parser
disagreement enters a separate queue rather than being resolved in favor of the
system under evaluation. Malicious samples that cannot be redistributed remain
identified by hashes and provenance in the isolated environment.

The fourth threat is dependency and platform variation. Codec libraries,
operating-system process controls, filesystem semantics, and compiler versions
can change behavior. Linux-specific address-space, CPU, file-size, process-group,
and no-new-privileges controls are not claimed on platforms where the
implementation does not apply them. The run manifest pins each dependency and
the target triple. Repeating the experiment on another platform is a separate
external-validity study.

The fifth threat is shared implementation. The reconstruction handler and
independent validator are separate modules, but they can still share libraries,
developers, and interpretations. External differential tools reduce but cannot
eliminate common-mode failure. Standards themselves can permit alternatives or
contain ambiguities. The accepted language is therefore the explicit Aura
profile, not an assertion that every standards-conforming consumer will behave
identically.

The sixth threat concerns availability. Fail-closed behavior can be abused to
deny service by submitting expensive or unsupported objects. Resource limits,
cancellation, tenant policy, and rate controls bound individual work, but a
complete distributed denial-of-service analysis lies outside this manuscript.
Deployment must combine per-object bounds with admission control and endpoint
resource governance.

\subsection{Research directions}

The nearest empirical task is execution of the frozen server protocol and
publication of the resulting manifests, raw observations, analysis scripts,
and minimized failures. Any release blocker discovered by the campaign must be
resolved or reported before numerical claims are promoted into the manuscript.

Two evidence-package synchronization items precede that freeze.  The
machine-readable alias register inspected at the pinned implementation commit
does not yet list \texttt{audio/m4a}, although the runtime maps it to the M4A
audio profile.  The fidelity adapter also recognizes APNG and GIF profile names
as animated but does not yet include
\texttt{webp-animated-canonical-v1}.  These omissions do not change the runtime
reconstruction result, but they prevent the current evidence package from
supporting alias-parity and animated-WebP fidelity claims.  The scientific
release predicate must remain false until both registries are synchronized and
their regression checks pass.

The nearest implementation task is to extend semantic coverage without
weakening profile precision. APNG, animated WebP, retained audio bitstreams,
and constrained video profiles require semantic paths before they can support
codec-payload neutralization claims.  Animated WebP now has a complete
structural reconstruction and validation profile, including a bounded
post-reconstruction frame decode, but it still retains compressed frame
payloads. Fragmented
ISO BMFF and broader Matroska features require separate grammars and offset
arguments rather than permissive extension of the current validators.

A deeper formal direction is machine-checked correspondence between selected
parser state machines, canonical serializers, and the acceptance predicate.
This work would need to separate verified native parsers from unverified codec
semantics and foreign processes. Another direction is authenticated policy
attestation between messaging endpoints. Such an extension could allow a sender
to know which reconstruction profile a receiver claims to enforce, but it would
change the messaging protocol and introduce new downgrade and privacy questions.

Finally, semantic safety mechanisms can be composed above canonical media
without conflating their evidence. Perceptual classifiers, provenance
credentials, child-safety policies, phishing analysis, and steganalysis can
consume only clean canonical outputs while retaining their own verdicts,
datasets, and uncertainty. The canonical boundary then serves as a controlled
input layer rather than a universal content judgment.
